\section{讨论与限制}

本文的主要结果不是简单证明 Transformer 优于 GRU，而是说明多时域区域航空预测需要两层时域适配。第一层是多尺度 horizon-query 编码：不同预测提前量应分别读取细尺度和粗尺度历史证据。第二层是 Horizon-Router：当短时局部机动、粗尺度稳定趋势和完整历史多尺度趋势同时存在时，时域级 router 可以为不同 horizon 选择更合适的 expert 组合。

区域化表述本身也很重要。在开放运行数据上，尤其是 60 到 120 分钟预测时域，长距离点式轨迹回归容易制造并不可靠的精确感。H3 区域概率把不确定性显式化，使 top-$k$ 排序、NLL、Brier 和 ECE 能够在同一离散概率空间中比较。本文本轮主结果刻意不报告预测集合面积，以免把额外集合构造与模型概率质量混在一起；预测集合大小、空间面积和覆盖风险之间的关系需要在后续用更独立的协议评价。

当前论文也有明确边界。首先，虽然 H3 表示在空间上是通用的，但本文实验使用的是 ADRR 中一个月商业航班运行档案，不应被解读为已经覆盖所有航空运行场景。更广泛的跨区域、跨季节和跨数据源验证仍然需要后续实验。

其次，当前模型保留 gap indicator，但没有显式建模缺失观测造成的不确定性。模型也不是 GRU-D、连续时间潜变量动力学或 conformal regional prediction 意义上的缺失感知预测器。这一部分应作为下一阶段工作，而不是当前论文的隐含贡献。

第三，路由训练机制仍需要完整规模逐项消融。当前 full-scale 证据主要体现在最终 Horizon-Router 相较 Multiscale Transformer 的提升；若要把 horizon-wise distillation、router supervision、load balancing 和 expert auxiliary loss 分别作为强因果贡献，还需要在相同数据划分和训练预算下补齐 router only、router + horizon-wise distillation、router + horizon-wise distillation + loss-aware routing regularization 三组实验。

最后，外部航空基线受任务不匹配限制较大。现有许多文献处理短时坐标预测、终端区交互，或依赖当前 ADRR 派生 artifact 中不可用的密集上下文输入。本文因此采用适配后的 TrAISformer-H3 和 FlightBERT++-H3 作为强基线，并明确固定数据、输出头和评价协议。该选择保证了同表比较的可解释性，但也意味着本文不是对原论文任务的逐字复现。
